\chapter{Writing a Task-Safe Peripheral Driver}
\label{ch:driver-design}\index{driver design}\index{task-safe driver}

Chapters~\ref{ch:hal}--\ref{ch:hal3} \emph{used} the drivers. This chapter shows
how to \emph{write} one, and -- just as important -- why it is built the way it
is. The shape described here is a project rule: every driver in the HAL for a
peripheral that is \textbf{exclusive} (cannot be shared between tasks) and can be
\textbf{reconfigured after setup} follows it exactly. By the end you should be
able to add a new one and get the same guarantees for free.

\section{The problem}

Consider a typical controller -- a SPI host, a UART, a CAN controller. Three
facts make it dangerous to share:

\begin{enumerate}
\item \textbf{It is one physical resource.} Two tasks driving the same FIFO or
  the same configuration registers at once corrupt each other's transfers.
\item \textbf{It is configured in stages.} You bring it up once (clocks, frame
  format), then later you may re-route pins, change the baud rate, flip a line's
  polarity -- often interleaved with real traffic.
\item \textbf{The registers are global.} On a bare-metal target the peripheral is
  a fixed memory-mapped block. Any code that can name it can poke it.
\end{enumerate}

A naive driver exports \texttt{Configure}, \texttt{Write}, \texttt{Read} as plain
procedures over a global. Nothing stops task B from calling \texttt{Set\_Baud}
while task A is mid-\texttt{Write}, or from starting a transfer it has no business
starting. The corruption is intermittent and miserable to debug. We want a design
where those mistakes are \emph{impossible to write}, not merely discouraged.

\section{The ownership model}\index{ownership gateway}\index{session}\index{capability}

The cure is to make \emph{ownership} a value you hold, and to make every operation
that touches the live controller require it.

\begin{description}[style=nextline]
\item[A \texttt{Session} is a capability.] \texttt{Acquire} hands a task a
  \texttt{Session} -- a limited, controlled handle that represents exclusive
  ownership of the controller. ``Limited'' means it cannot be copied (two tasks
  can never share one), and ``controlled'' means it releases automatically when it
  leaves scope, even on an exception. While you hold it, no other task can
  acquire the controller.
\item[Three things become impossible.] You cannot acquire a controller that was
  never set up (\texttt{Not\_Initialized}); you cannot transfer or reconfigure
  without holding it (\texttt{Not\_Owned}); and you cannot leak the lock, because
  the handle's finalizer always releases.
\end{description}

The \texttt{Session} is the public face. The interesting part is the machinery
behind it that makes ``you cannot reconfigure without holding it'' a structural
fact rather than a rule a future maintainer must remember.

\section{Three layers}\index{Engine pattern}\index{private child}

A compliant driver is three packages:

\begin{enumerate}
\item \textbf{The Engine} (\texttt{ESP32S3.<Drv>.Engine}, a \emph{private child}).
  It owns \emph{all} register access and exposes a \texttt{Bus} handle. It is the
  only unit in the whole program that names the peripheral's registers.
\item \textbf{The State gateway} (a package nested in the parent's \emph{body}). It
  holds the one configured \texttt{Bus} handle and hands it out through a single
  function, \texttt{Owned (S)}, which raises unless \texttt{S} holds the
  controller.
\item \textbf{The public package} (\texttt{ESP32S3.<Drv>}). It declares the
  \texttt{Session}, the exceptions, \texttt{Setup}, \texttt{Acquire}, and the
  operations -- each of which is a one-liner that reaches the hardware only
  through \texttt{Owned}.
\end{enumerate}

The rest of this chapter builds a small driver, \texttt{ESP32S3.Widget}, one layer
at a time.

\begin{figure}[ht]
\centering
\resizebox{\linewidth}{!}{%
\begin{tikzpicture}[node distance=12mm]
\node[blk, text width=2.4cm] (app) {Application\\task};
\node[blk, right=of app, fill=blue!8, text width=3.4cm] (pub)
  {Public package\\\scriptsize limited \texttt{Session} +\\\scriptsize protected \texttt{Guard}};
\node[blk, right=of pub, fill=blue!8, text width=2.8cm] (gw)
  {\texttt{Owned (S)}\\\scriptsize gateway\\\scriptsize (raises if not held)};
\node[blk, right=of gw, fill=orange!10, text width=2.6cm] (eng)
  {\texttt{Engine}\\\scriptsize private child};
\node[blk, right=of eng, fill=gray!12, text width=2.0cm] (reg) {Registers};
\draw[flow] (app) -- node[above,font=\scriptsize]{Acquire} (pub);
\draw[flow] (pub) -- (gw);
\draw[flow] (gw)  -- (eng);
\draw[flow] (eng) -- (reg);
\end{tikzpicture}}
\caption{The task-safe driver skeleton: a transfer reaches the registers only via
the ownership-checked \texttt{Owned} gateway, and only the private \texttt{Engine}
names them. Forgetting the check is structurally impossible.}
\label{fig:driver-skeleton}
\end{figure}

\section{Layer 1 -- the Engine (the only code that sees the registers)}

The engine is a \emph{private child}: a unit named \texttt{ESP32S3.Widget.Engine}
can be \texttt{with}-ed only from inside the \texttt{ESP32S3.Widget} subtree.
Application code -- and even other drivers -- physically cannot reach it.

\begin{lstlisting}
private package ESP32S3.Widget.Engine is
   type Bus is private;                       --  an opaque, configured controller

   function  Open (Clock_Hz : Positive) return Bus;
   procedure Configure_Pins (B : Bus; Tx, Rx : ESP32S3.GPIO.Optional_Pin);
   procedure Transfer (B : Bus; Tx, Rx : System.Address; Length : Natural);
private
   type Bus is record
      Regs  : Periph_Ref := null;             --  pointer to the register block
      Valid : Boolean    := False;
   end record;
end ESP32S3.Widget.Engine;
\end{lstlisting}

The engine \emph{body} is the one place that \texttt{with}s the SVD register
package. It does the bit-twiddling -- the same code you would have written
anyway -- but now it is sealed behind the \texttt{Bus} handle. \texttt{Open}
brings the controller up and returns a \texttt{Bus}; every other operation takes a
\texttt{Bus} and works through it. Crucially, \textbf{the parent package never
\texttt{with}s the SVD registers at all} -- so once the engine exists, the global
controller is unreachable from anywhere except through a \texttt{Bus} handle.

\section{Layer 2 -- ownership: the Session and the Guard}\index{protected object}\index{limited type}\index{controlled type}

The public package declares the handle and a per-controller lock. The
\texttt{Session} is limited and controlled:

\begin{lstlisting}
package ESP32S3.Widget is
   type Session is limited private;

   Not_Initialized : exception;   --  Acquire before Setup
   Not_Owned       : exception;   --  use without holding the controller

   procedure Setup (Clock_Hz : Positive := 1_000_000);          --  startup only
   procedure Acquire (S : in out Session);
   procedure Configure_Pins (S : Session; Tx, Rx : ESP32S3.GPIO.Optional_Pin);
   procedure Transfer (S : Session; Tx, Rx : System.Address; Length : Natural);
   procedure Release (S : in out Session);
private
   type Session is new Ada.Finalization.Limited_Controlled with record
      Active : Boolean := False;
   end record;
   overriding procedure Finalize (S : in out Session);          --  auto-release
end ESP32S3.Widget;
\end{lstlisting}

In the body, a protected object arbitrates ownership. Its guarded section is
tiny -- it only flips a flag -- so a task is never blocked inside it while a
slow transfer runs:

\begin{lstlisting}
protected Guard is
   entry    Acquire;
   procedure Release;
private
   Held : Boolean := False;
end Guard;

protected body Guard is
   entry Acquire when not Held is begin Held := True;  end Acquire;
   procedure Release         is begin Held := False; end Release;
end Guard;
\end{lstlisting}

\texttt{entry Acquire when not Held} is the whole mutual-exclusion story: a second
task calling \texttt{Acquire} \emph{suspends} on the barrier until the first
releases. (For a peripheral with several identical instances -- e.g. two SPI
hosts -- this is an \texttt{array (Host) of Guard} and a matching array of
handles; the singleton case is shown here for clarity.)

\section{Layer 3 -- the gateway that cannot be bypassed}\index{Owned gateway}

Here is the heart of the design. The configured \texttt{Bus} lives inside a
package nested in the \emph{body}, so no other part of the driver can name it. The
only way out is \texttt{Owned}:

\begin{lstlisting}
package State is
   procedure Open (Clock_Hz : Positive);       --  the Setup path
   function  Ready return Boolean;
   function  Owned (S : Session) return E.Bus; --  THE gateway
end State;

package body State is
   The_Bus : E.Bus;                            --  hidden here; unreachable elsewhere
   Is_Set  : Boolean := False;

   procedure Open (Clock_Hz : Positive) is
   begin
      The_Bus := E.Open (Clock_Hz);
      Is_Set  := True;
   end Open;

   function Ready return Boolean is (Is_Set);

   function Owned (S : Session) return E.Bus is
   begin
      if not S.Active then
         raise Not_Owned with "Widget used without holding it -- Acquire first";
      end if;
      return The_Bus;
   end Owned;
end State;
\end{lstlisting}

Now every public operation is a single line that funnels through \texttt{Owned},
and the ownership check happens in exactly one place:

\begin{lstlisting}
procedure Setup (Clock_Hz : Positive := 1_000_000) is
begin
   State.Open (Clock_Hz);
end Setup;

procedure Acquire (S : in out Session) is
begin
   if not State.Ready then
      raise Not_Initialized with "Widget acquired before Setup";
   end if;
   Guard.Acquire;                              --  suspends until free
   S.Active := True;
end Acquire;

procedure Configure_Pins (S : Session; Tx, Rx : ESP32S3.GPIO.Optional_Pin) is
begin
   E.Configure_Pins (State.Owned (S), Tx, Rx); --  Owned raises unless held
end Configure_Pins;

procedure Transfer (S : Session; Tx, Rx : System.Address; Length : Natural) is
begin
   E.Transfer (State.Owned (S), Tx, Rx, Length);
end Transfer;

procedure Release (S : in out Session) is
begin
   if S.Active then
      S.Active := False;
      Guard.Release;
   end if;
end Release;

overriding procedure Finalize (S : in out Session) is
begin
   Release (S);                                --  scope exit / exception unwind
end Finalize;
\end{lstlisting}

Look at \texttt{Configure\_Pins} and \texttt{Transfer}: there is \emph{no}
\texttt{if S.Active then} in either. They cannot reach the hardware without
calling \texttt{State.Owned (S)}, and \texttt{Owned} is the one place that checks.
A maintainer adding a new operation tomorrow has no choice but to go through the
same door -- there is no \texttt{Bus} in scope to grab any other way.

\section{The limited-handle variant}

If the engine's \texttt{Bus} is a \emph{limited} type -- typically because it
holds something limited, like a finalization-controlled DMA channel -- it cannot
be returned by value. Make the hidden instance \texttt{aliased} and have
\texttt{Owned} return an access to it:

\begin{lstlisting}
The_Bus : aliased E.Bus;
function Owned (S : Session) return access E.Bus is
begin
   if not S.Active then raise Not_Owned with "..."; end if;
   return The_Bus'Access;
end Owned;
--  callers: E.Transfer (State.Owned (S).all, Tx, Rx, Length);
\end{lstlisting}

In the HAL, SPI and I2S (whose engine \texttt{Bus} is limited) and LCD (whose
handle owns a GDMA channel) use this form; I2C, UART and TWAI return the handle by
value. Nothing else about the pattern changes.

\section{Why it is built this way}\index{RAII}

Each layer earns its place by removing a specific class of mistake:

\begin{description}[style=nextline]
\item[Private-child engine $\rightarrow$ the registers have one owner.] Because the
  parent never \texttt{with}s the SVD package, the controller global is
  unreachable outside the engine. The ``any code can poke it'' hazard is gone at
  the visibility level, not by convention.
\item[Hidden handle + single \texttt{Owned} gateway $\rightarrow$ the guard cannot be
  forgotten.] The earlier, weaker approach put an \texttt{if S.Active} check at the
  top of each operation. That works until someone adds an operation and forgets
  the check. Hiding the only \texttt{Bus} behind \texttt{Owned} makes the check
  structural: forgetting it is not a bug you can write, because there is no other
  way to obtain a \texttt{Bus}.
\item[Limited \texttt{Session} $\rightarrow$ ownership cannot be duplicated.] A
  non-copyable handle means two tasks can never hold the same controller, and a
  handle can never outlive its scope as a dangling copy.
\item[Controlled \texttt{Session} (RAII) $\rightarrow$ the lock cannot leak.] The
  finalizer releases on every exit path, including exception unwinding. A fault
  between \texttt{Acquire} and \texttt{Release} hands the controller back instead
  of deadlocking the next task forever.
\item[Tiny protected section $\rightarrow$ ownership without serialising work.] The
  lock is held only long enough to flip a flag; the blocking transfer runs
  outside it. Tasks contend for \emph{ownership}, never for the bus itself.
\item[Two named exceptions $\rightarrow$ misuse fails loudly.] \texttt{Not\_Initialized}
  and \texttt{Not\_Owned} turn ``used in the wrong state'' into an immediate,
  explanatory error instead of a silent no-op or a corrupted transfer.
\end{description}

\section{The benefits}

\begin{itemize}
\item \textbf{Mistakes become unrepresentable.} ``Configure a port you don't
  own,'' ``transmit from the wrong task,'' ``poke the registers directly'' --
  none of these can be written. The compiler and the visibility rules enforce
  the discipline, so code review does not have to.
\item \textbf{One pattern, every peripheral.} Because all six \texttt{Session}
  drivers share this skeleton, learning one teaches all of them, and a new driver
  is a fill-in-the-blanks exercise rather than a fresh concurrency design.
\item \textbf{The unsafe core is testable in isolation.} The engine is plain,
  lock-free register logic; the parent is plain ownership logic. Each can be
  reasoned about -- and the engine bench-tested -- on its own.
\item \textbf{RAII hygiene for free.} Scope-based release means the common
  ``acquire, do work, release'' path needs no explicit \texttt{Release} and is
  exception-safe by construction.
\item \textbf{It costs almost nothing.} The handle is a couple of words, the
  gateway is one comparison, and the protected section is a flag flip. The safety
  is structural, not runtime overhead.
\end{itemize}

\section{When the rule applies (and when it does not)}

Use this pattern for a driver that is \textbf{both} exclusive (cannot be shared
between tasks) \textbf{and} reconfigurable after setup. That is the combination
that makes uncoordinated access dangerous.

It does \emph{not} apply to lock-free or shareable peripherals. \texttt{GPIO}'s
set/clear/read are hardware-atomic; \texttt{RNG} is a stateless read; the
\texttt{Temperature}, \texttt{RTC} and crypto blocks either are internally
serialised by a small protected object or have nothing a second task can corrupt.
Those drivers stay simple on purpose and run under every runtime profile,
including light-tasking, where the controlled \texttt{Session} types (which need
finalization) are not available. Reach for the full ownership machinery only when
the hazard is real; for everything else, the simplest safe thing is better.

\section{Letting the application plug in: dispatch or callback}
\label{driver:plugin}\index{callback}\index{dispatching}\index{closure-free}

A driver often needs the \emph{application} to supply a piece of behaviour: assert
a chip-select line, pin a fresh socket to one interface, decide what a Modbus
register read returns. In most languages the reach-for tool is a \emph{closure}.

A closure is a subprogram bundled with the variables it uses but does not itself
declare -- the surrounding state it ``closes over.'' A nested subprogram that reads
a variable from its enclosing scope is the everyday case, and on a driver boundary
it shows up like this:

\begin{lstlisting}
procedure Configure (Bus : in out E.Bus; CS : ESP32S3.GPIO.Pin_Id) is
   procedure Select_It (Active : Boolean) is
   begin
      ESP32S3.GPIO.Write (CS, Active);     -- CS is captured from Configure
   end Select_It;
begin
   Open (Bus, Chip_Select => Select_It'Access);   -- hand the nested proc out
end Configure;
\end{lstlisting}

\texttt{Select\_It} declares only \texttt{Active} but uses \texttt{CS}, so to be
callable later it has to remember \emph{which} \texttt{CS} -- the one from the
\texttt{Configure} call that produced it. That pairing of code with captured state
is the closure. A subprogram that captures \emph{nothing} -- using only its own
parameters, globals and library-level state -- is not a closure but a plain
function pointer, and that distinction is the whole of what follows.

A closure is exactly what this platform forbids. A subprogram that captures its
environment compiles to a \emph{trampoline} -- a scrap of executable code GCC
writes onto the stack, which is where the captured \texttt{CS} gets bound -- and on
the ESP32-S3 the stack lives in non-executable data memory, so calling such a
trampoline faults outright. The HAL is built under
\texttt{pragma Restrictions (No\_Implicit\_Dynamic\_Code)}, which turns that into a
\emph{compile-time} refusal rather than a runtime crash (the same fault, seen from
the debugger's side, is the war story in Chapter~\ref{ch:debug}).

So a driver's extension point must need \emph{no} captured state. Two shapes
satisfy that, and which one a driver picks comes down to how much per-instance
state the hook carries.

\textbf{A closure-free callback with an explicit context.} The hook is a
library-level access-to-subprogram, and any state it needs travels through an
opaque context parameter the driver hands back to it -- not captured, passed. The
SPI driver's application chip-select is this shape (Chapter~\ref{ch:hal}):

\begin{lstlisting}
type CS_Select is access procedure (Ctx : System.Address; Active : Boolean);
--  the app's procedure is library-level; per-device state rides in Ctx
\end{lstlisting}

It is light and runs under every profile. The cost is that \texttt{Ctx} is
untyped -- the callback converts it back itself -- so it suits a hook that is
essentially stateless or one-shot: ``drive this pin,'' ``call \texttt{Set\_Interface}
on this socket.''

\textbf{Dispatching through a tagged type.} You derive a type, keep your data in
its fields, and override the operations you support. The Modbus slave is this
shape (Chapter~\ref{ch:modbus}):

\begin{lstlisting}
type Device is new Modbus.Slave.Server with record
   Holding : Word_Array (0 .. 63) := (others => 0);   -- your storage, in the object
end record;
overriding procedure On_Read_Holding_Registers (Self : in out Device; ...);
\end{lstlisting}

Here \texttt{Self} \emph{is} the context: no \texttt{Ctx} to thread, no conversion,
and the call dispatches through a static vtable -- a table, not generated code, so
\texttt{No\_Implicit\_Dynamic\_Code} is satisfied. The set of primitives you
override expresses exactly which operations the device implements; the ones you
leave alone take the base type's default (for Modbus, an \texttt{Illegal\_Function}
reply). This suits a hook whose context is rich and structured, or a \emph{family}
of related operations with sensible defaults.

The rule of thumb: a stateless or one-shot hook wants the closure-free callback
(carry any scrap of context in an opaque pointer); a hook with rich per-instance
state, or a group of related operations, wants a tagged type you derive and
override. Both exist for one reason -- to recover the convenience of a capturing
closure without the trampoline this chip cannot run.
